%%%%%%%%%%%%%%%%%%%%%%%%%%%%%%%%%%%%%%%%%
% Dreuw & Deselaer's Poster
% LaTeX Template
% Version 2.0 (February 18, 2023)
%
% This template originates from:
% https://www.LaTeXTemplates.com
%
% Authors:
% Vel (vel@latextemplates.com)
% Philippe Dreuw and Thomas Deselaers (https://github.com/deselaers/latex-beamerposter)
%
% License:
% CC BY-NC-SA 4.0 (https://creativecommons.org/licenses/by-nc-sa/4.0/)
%
% NOTE: The bibliography needs to be compiled using the bibtex engine.
%
%%%%%%%%%%%%%%%%%%%%%%%%%%%%%%%%%%%%%%%%%

%----------------------------------------------------------------------------------------
%	PACKAGES AND OTHER DOCUMENT CONFIGURATIONS
%----------------------------------------------------------------------------------------

\documentclass{beamer} % Use the beamer base class

\usepackage[
	orientation=portrait, % Portrait orientation
	size=a1, % Paper size
	scale=1, % Scale, it's important to adjust this so your content fits nicely in the template
]{beamerposter} % Use the beamerposter package to create the layout

\usetheme{I6pd2} % Use the I6pd2 theme supplied with this template

\usepackage{changepage} % Required for temporarily indenting text blocks

\usepackage{amsmath,amsthm,amssymb,latexsym} % For including math equations, theorems, symbols, etc
\usepackage{physics}
\usepackage{gfsdidot} % Use the GFS Didot font

\usepackage{booktabs} % Top and bottom rules for tables

\usefonttheme{professionalfonts}
%\usefonttheme[onlymath]{structureitalicserif} % Uses Computer Modern for maths

\graphicspath{{Figures/}} % Location of figure images

%----------------------------------------------------------------------------------------
%	TITLE SECTION
%----------------------------------------------------------------------------------------

\title{\Huge Can disease spread be modelled by a reaction-diffusion equation?} % Poster title

\author{Oliver Clift\textsuperscript{1}, James Emberton\textsuperscript{1}, Isaac Millen\textsuperscript{1}, Nicholas Mitchiner\textsuperscript{1}, Edward Phillips\textsuperscript{1} and Dominik Szablonski\textsuperscript{1}} % Author(s)

\institute{\textsuperscript{1}Department of Physics, University of Manchester} % Institution(s)

%----------------------------------------------------------------------------------------
%	FOOTER TEXT
%----------------------------------------------------------------------------------------

\newcommand{\leftfoot}{} % Left footer text

\newcommand{\rightfoot}{} % Right footer text

%\setbeamertemplate{footline}{} % Uncomment this line to hide the footer

%----------------------------------------------------------------------------------------

\begin{document}
	
\begin{frame}[t] % The whole poster is enclosed in one beamer frame, the [t] parameter aligns everything to the top

\begin{columns}[t] % Begin multi-column layout, the [t] parameter aligns each column's content to the top

\begin{column}{0.02\textwidth}\end{column} % Empty column for horizontal whitespace

\end{columns}

\hspace{0.015\textwidth}%
\begin{minipage}{0.97\textwidth}
\begin{block}{Abstract}
	In this study, we modelled the spread of disease in Manchester through a reaction-diffusion equation, comparing this to COVID-19 data in Boston. The main assumptions of the model were a constant population density, a recovery time of 7 days, and negligible effects on immunocompromised individuals. The Euler method was used to solve the equation computationally. Our model successfully recreated spatial disease spread when compared to COVID-19 data in Boston, although it diverged from the data after roughly 30 days. The model could be improved by accounting for more advanced assumptions including a travelling population, multiple outbreak centres and behavioural changes.

\end{block}
\end{minipage}

\begin{columns}[t]
\begin{column}{0.02\textwidth}\end{column}

\begin{column}{0.465\textwidth} % Start the first content column


%----------------------------------------------------------------------------------------
%	INTRODUCTION
%----------------------------------------------------------------------------------------
            
\begin{block}{Introduction}
	Modelling disease spread is important for governments have appropriate plans in case of large-scale pandemics. The recent Covid pandemic was not handled very well (give some examples, probably cite some things here). With this motivation, we attempted to model the disease spread within greater Manchester. Using our knowledge of physics, we supposed that the spread of a disease like Covid could be modelled by a diffusion equation.
	
	\bigskip
	
	However, due to mass conservation, a standard diffusion equation is incapable of accurately modelling disease spread, because diseases are inherently create new infections. Therefore, for disease modelling, a reaction-diffusion equation is much so that the number of infections at any time is not conserved. This allows us to model new infections and patient recovery.
\end{block}

%----------------------------------------------------------------------------------------
%	MATERIALS
%----------------------------------------------------------------------------------------

\begin{block}{Theory - Reaction-Diffusion Equation} 
	We attempt to model the spread of disease via a reaction-diffusion equation,
	\begin{equation}
		\pdv{I}{t} = D\grad^2I + \beta(x,y)SI 	- \gamma I,
	\end{equation}
	which models spatial densities $\rho$, which are parametrised such that,
	\begin{equation}
		\rho = S + I + R,
	\end{equation}
	where,
	\begin{itemize}
		\item $S(x,y,t)$: Susceptible population.
		\begin{itemize}
			\item Those who can catch disease.
		\end{itemize}
		\item $I(x,y,t)$: Infected population.
		\begin{itemize}
			\item Those who have the disease.
		\end{itemize}
		\item $R(x,y,t)$: Recovered population.
		\begin{itemize}
			\item Those who are no longer infected.
		\end{itemize}
	\end{itemize}
\end{block}

\begin{block}{Results: Simulation vs Data}
	Infection of COVID-19 was modelled in Greater Manchester with an initial transmissibility of $\beta_0 = 0.48$~per day taking into account weather effects \cite{TALIB2021}
	while using a diffusion coefficient of $D=0.85$~m$^{2}$~s$^{-1}$, and a mean recovery time of $\tau = 7$~days \cite{NHS2023}
	and using the spatial density given for each county of Greater Manchester from the March 2021 Census data \cite{CENSUS2021}.
	The model was compared to the COVID-19 infection map of Boston \cite{KUMAR2021} as it is a city of comparable size thus allowing us to 
	clearly compare the accuracy of our model. As seen our model is comparable to the Boston infection map up to the first 30 days.
	After that, the model diverges and while spreading out the infection at the initial points in the centre of Greater Manchester
	disappears as opposed to the Boston heat map.
\end{block}

\begin{block}{Results: Model Figure}
	\begin{figure}
		\includegraphics[width=\linewidth]{infection_spread_model_v3.png}
		\caption{COVID-19 infection over time in Greater Manchester up to 100 days.
			Using parameters of initial transmission rate $\beta_0 = 0.48$~per day and diffusion coefficient of $D = 0.85$~m$^2$~s$^{-1}$ 
			and a mean recovery time of $\tau = 7$~days.}
		\label{model}
	\end{figure}
\end{block}
%----------------------------------------------------------------------------------------

\end{column} % End of the first column

\begin{column}{0.03\textwidth}\end{column} % Empty column for horizontal whitespace
 
\begin{column}{0.465\textwidth} % Start the second content column

%----------------------------------------------------------------------------------------
%	RESULTS
%----------------------------------------------------------------------------------------

%------------------------------------------------

\begin{block}{Results: Spread of Covid From Data}
	\begin{figure}
		\includegraphics[width=\linewidth]{Bostton_auto_sprawl_figure.png}
		\caption{Covid-19 infection "Heatmap of infected (exposed and infectious) individuals every seven days: (a) Auto Sprawl;"\cite{KUMAR2021} of Boston over 84 days.}
\label{Boston_data}
	\end{figure}
\end{block}

\begin{block}{Discussion}
We found that though our model was successful at modelling the initial spread of the Covid 19 pandemic. In particular, the initial growth seen in Fig.\ref{model} matches near perfectly data seen in Fig. \ref{Boston_data}. However, after this, the model starts to diverge from real life data. This is likely because we made the assumption that people’s behaviour wouldn’t change when infected and the nature of the Reaction-Diffusion equation used in the model. In reality, as time went on, and more was understood about Covid, people started to self isolate when they first experienced symptoms, and the first lockdowns started to occur in March 2020 \cite{iog}. However, this doesn’t make the idea of using diffusion-reaction useless or invalid in all cases. Instead, it implies that the best use of such a model is to predict what would happen should the spread of the virus not be contained for the initial spread. 

\bigskip

Our model also fails to account for individuals travelling long distances, instead assuming that everyone remains in roughly the same location. In reality, this isn’t true for multiple reasons. Firstly, as a big city, people commute into central Manchester from the likes of Salford, Stockport, and Tameside both to shop and work. These people are likely to get infected in Manchester and travel back to their homes, causing the pandemic to spread to outer regions sooner than in the model. Greater Manchester isn’t closed either with both national and international travel happening in and out daily. Whilst an infected person leaving Manchester wouldn’t change much, at least outside of the very early days, an infected person entering Manchester could lead to a new infection centre. In reality, it turns out that, “connectivity matters more than density in the spread of COVID-19” \cite{HAMIDI2020102378}.

\bigskip

A diffusion-reaction model also fails to account for the virus, “Sticking around”. In our model, after around 30 days the area where our “patient zero” first appeared had been near completely cleared out as everyone in the area had caught the illness and recovered. In reality, not everyone in an area gets ill immediately, so the illness will stay around for longer rather than just spread to other areas. Despite that, the idea that a virus may exhaust all of its potential inhabitants in an area does have historical backing, though this is more common without lockdown protocol.
PARAGRAPH NEEDS SOURCING AND NEEDS DATA.

\end{block}

%----------------------------------------------------------------------------------------
%	CONCLUSIONS
%----------------------------------------------------------------------------------------

\begin{block}{Conclusions} 
	In conclusion, the reaction-diffusion equation successfully captured the initial spread of COVID-19, but diverged from real-world data at later times. This is due to simplifying assumptions, such as a lack of behavioural changes, a travelling population, and the virus not lingering in certain areas, which a better model for disease spread would account for.
\end{block}
%----------------------------------------------------------------------------------------
%	REFERENCES
%----------------------------------------------------------------------------------------

\begin{block}{References}
	%\nocite{*} % Insert publications even if they are not cited in the poster
	\small % Reduce font size
	\vspace{-1ex} % Pull up slightly
	\bibliographystyle{unsrt} % Bibliography style
	\bibliography{sample.bib} % Bibliography file
\end{block}


%----------------------------------------------------------------------------------------

\end{column} % End of the second column

\begin{column}{0.02\textwidth}\end{column} % Empty column for horizontal whitespace

\end{columns} % End of all the columns in the poster

\end{frame} % End of the enclosing frame

%----------------------------------------------------------------------------------------

\end{document}
